\section*{अध्याय \ref{ch:b3:molecular-evolution} --- आणविक विकास और जातिवृत्तजीनोमिकी}

\begin{solution}{exo:b3:molecular-evolution:1}
$k = \mu$: प्रति स्थल प्रति पीढ़ी उदासीन \omterm{thm:b3:molecular-evolution:neutral}{प्रतिस्थापन दर} उत्परिवर्तन-दर
के बराबर है। हर पीढ़ी $2N\mu$ नए उदासीन उत्परिवर्तन उठते हैं और हर एक
के स्थिर होने की प्रायिकता $1/2N$ है; कोई बड़ी समष्टि अधिक उत्परिवर्तन
बनाती है और हर एक को अनुपात में छोटी संभावना देती है, और ये दोनों असर
ठीक-ठीक कट जाते हैं।
\end{solution}

\begin{solution}{exo:b3:molecular-evolution:2}
\omterm{def:b3:genomics:comparative}{ऑर्थोलॉग}: दो जातियों के वे जीन जो उनके साझा पूर्वज के किसी एक जीन से
उतरे हों --- मानव $\beta$-ग्लोबिन और चूहे का $\beta$-ग्लोबिन।
\omterm{def:b3:genomics:comparative}{पैरालॉग}: किसी एक जीनोम के वे जीन जो किसी द्विगुणन से उतरे हों ---
मानव $\alpha$- और $\beta$-ग्लोबिन, या $\beta$ तथा $\gamma$। \omterm{def:b3:molecular-evolution:duplication}{आभासी जीन}:
कोई ऐसी क्षयग्रस्त प्रति जो अब कोई प्रोटीन कूटबद्ध नहीं करती --- मानव
$\psi\beta$ जीन, जो $\beta$-ग्लोबिन गुच्छे में है, विराम-कूटकों समेत।
\end{solution}

\begin{solution}{exo:b3:molecular-evolution:3}
$\omega$ ऐमीनो-अम्ल बदलती प्रतिस्थापन-दर की तुलना मौन प्रतिस्थापन-दर से
करता है, जो वह उदासीन आधार-रेखा है। $0.02$: प्रबल \omterm{def:b3:molecular-evolution:dnds}{शोधक वरण}, यानी
ऐमीनो-अम्ल बदलावों का \qty{98}{\%} हटा दिया गया --- ऐक्टिन, \omterm{def:b3:chromatin-epigenetics:nucleosome}{हिस्टोन}
H3। $1.0$: कोई बंधन नहीं --- कोई \omterm{def:b3:molecular-evolution:duplication}{आभासी जीन}। $2.5$: \omterm{def:b3:molecular-evolution:dnds}{धनात्मक वरण}, यानी
ऐमीनो-अम्ल बदलावों का पक्ष लिया गया --- MHC की पेप्टाइड-बंधक खाँच,
HIV का \omterm{def:b3:virology:virus}{परिवेष्टन}, रोमंथी लाइसोज़ाइम।
\end{solution}

\begin{solution}{exo:b3:molecular-evolution:4}
स्तनियों भर मिलाए गए प्रोटीन प्रति स्थल प्रति अरब वर्ष लगभग एक
प्रतिस्थापन से बदलते हैं; और अरब भर कूटलेखी स्थलों के किसी जीनोम तथा कुछ
वर्षों की किसी पीढ़ी भर यह प्रति पीढ़ी एक स्थिर हुए प्रतिस्थापन की, यानी
हर दो वर्ष में एक की, कोटि का है। हाल्डेन की वरण-लागत लगभग तीन सौ
पीढ़ियों में एक वरित प्रतिस्थापन की छूट देती है। वे दरें सौ या उससे
अधिक के किसी गुणक से भिन्न हैं, इसलिए लगभग हर स्थिर हुए बदलाव को उदासीन
होना ही चाहिए।
\end{solution}

\begin{solution}{exo:b3:molecular-evolution:5}
उदासीन: $1/2N = 10^{-4}$। $s = 0.001$, $4Ns = 20$: $(1 - \mathrm{e}^{-0.002})/
(1 - \mathrm{e}^{-20}) = 0.002$। $s = 0.01$: $0.0198$। $s = -0.001$: $(1 -
\mathrm{e}^{0.002})/(1 - \mathrm{e}^{20}) = 0.002\,\mathrm{e}^{-20} =
4\times 10^{-12}$। कोई भी प्रभावतः उदासीन नहीं: उसके लिए $|4Ns|
\lesssim 1$ चाहिए होता, यानी $|s| \lesssim 5\times 10^{-5}$; और प्रतिशत
के दसवें भाग का कोई वरण गुणांक पाँच हज़ार की किसी समष्टि में पहले ही
निर्णायक है।
\end{solution}

\begin{solution}{exo:b3:molecular-evolution:6}
बिना सुधार: $T = 0.12/(2\times 2\times 10^{-9}) = \qty{30}{Myr}$।
जूक्स--कैंटर: $d = -\tfrac{3}{4}\ln(1 - 0.16) = 0.131$, $T =
\qty{33}{Myr}$। वह सुधार दो का गुणक तब छूता है जब $-\tfrac{3}{4}
\ln(1 - 4p/3) = 2p$ हो, यानी $p \approx 0.6$ पर: तब तक वह कच्चा अंतर
अपनी अधिकांश सूचना खो चुका होता है।
\end{solution}

\begin{solution}{exo:b3:molecular-evolution:7}
$p_{N} = 7/700 = 0.010$, $p_{S} = 10/200 = 0.050$, $\omega = 0.2$:
यानी \omterm{def:b3:molecular-evolution:dnds}{शोधक वरण}, जो ऐमीनो-अम्ल बदलावों के चार-पाँचवें भाग हटा देता है।
$35$ तथा $10$ के साथ: $p_{N} = 0.05$, $\omega = 1$: कोई पकड़ में आता
बंधन नहीं --- यानी कोई \omterm{def:b3:molecular-evolution:duplication}{आभासी जीन}, या कोई ऐसा जीन जिसमें कुछ स्थलों पर
\omterm{def:b3:molecular-evolution:dnds}{धनात्मक वरण} दूसरों पर के \omterm{def:b3:molecular-evolution:dnds}{शोधक वरण} को संतुलित कर देता है।
\end{solution}

\begin{solution}{exo:b3:molecular-evolution:8}
$k = d/2T = 0.012/(1.3\times 10^{7}) = 9.2\times 10^{-10}$ प्रति स्थल
प्रति वर्ष; $\mu = 25k = 2.3\times 10^{-8}$ प्रति स्थल प्रति पीढ़ी;
$6\times
10^{9}\times 2.3\times 10^{-8} \approx 140$ नए उत्परिवर्तन प्रति
द्विगुणित जीनोम प्रति पीढ़ी। माता-पिता तथा संतानों के सीधे अनुक्रमण से
लगभग सत्तर मिलते हैं: वह अपसरण-आकलन उस बहुरूपता से बढ़ा हुआ है जो उस
पुरखा समष्टि में पहले से मौजूद थी, और जो उस दिखते विभाजन में कुछ लाख
पीढ़ियाँ जोड़ देती है।
\end{solution}

\begin{solution}{exo:b3:molecular-evolution:9}
$\mathrm{e}^{-T/2N} = \mathrm{e}^{-1} = 0.37$; असहमत अंश $\tfrac{2}{3}
\times 0.37 = 0.25$। \qty{10}{\%} के लिए: $\mathrm{e}^{-T/2N} = 0.15$, $T =
2N\ln(1/0.15) = \num{190000}$ पीढ़ियाँ।
\end{solution}

\begin{solution}{exo:b3:molecular-evolution:10}
$s \to 0$ पर: $1 - \mathrm{e}^{-2s} \approx 2s$ और $1 - \mathrm{e}^{-4Ns}
\approx 4Ns$, अनुपात $1/2N$। $4Ns \gg 1$ के लिए हर $1$ है और
अंश $\approx 2s$। $s < 0$ के लिए, जिसमें $4N|s| \gg 1$ हो: अंश $1 -
\mathrm{e}^{2|s|} \approx -2|s|$, हर $1 - \mathrm{e}^{4N|s|} \approx
-\mathrm{e}^{4N|s|}$, इसलिए $P \approx 2|s|\,\mathrm{e}^{-4N|s|}$।
$N = 10^{4}$ में उदासीन से सौ गुना नीचे: $2|s|\,\mathrm{e}^{-4\times 10^{4}|s|} =
5\times 10^{-7}$; और $|s| = 1.6\times 10^{-4}$ देता है $3.2\times 10^{-4}\times
\mathrm{e}^{-6.4} = 5.3\times 10^{-7}$। इसलिए $|s| \approx 1.6\times 10^{-4}$,
$4N|s| \approx 6.5$: यानी दस हज़ार व्यष्टियों में प्रतिशत के सौवें भाग
की कोई हानि ही स्थिरण को संयोग से सौ गुना बिरला बनाने के लिए काफ़ी है।
\end{solution}

\begin{solution}{exo:b3:molecular-evolution:11}
दूरी $a$ पर उस विभाजन-बिंदु को, दोनों रास्तों पर $O$ से, लेने पर, $a + m =
0.30$, $a + h = 0.24$, $m + h = 0.20$: $m - h = 0.06$, इसलिए $m = 0.13$, $h
= 0.07$, अनुपात $1.9$। कोई कड़ी प्रति-पीढ़ी घड़ी पीढ़ी-गिनतियों का
अनुपात देती, यानी $25/0.5 = 50$। देखा गया $1.9$ यह कहता है कि उस चूहे ने
पचास गुनी पीढ़ियों के बावजूद मनुष्य से प्रति वर्ष केवल दुगुना बदलाव जमा
किया: यानी प्रति पीढ़ी उत्परिवर्तन-दर मनुष्यों में कोई $25$ गुना ऊँची
होनी चाहिए --- प्रति पीढ़ी अधिक जनन-वंशीय कोशिका-विभाजन --- ताकि प्रति
वर्ष वे दोनों वंश-परंपराएँ केवल दो के किसी गुणक से भिन्न हों।
\end{solution}

\begin{solution}{exo:b3:molecular-evolution:12}
विराम बनाते प्रतिस्थापन $1000\times 2\times 10^{-9}\times
0.04 = 8\times 10^{-8}$ प्रति वर्ष से आते हैं: पहला लगभग
\qty{12}{Myr} बाद प्रत्याशित है (खाँचा-खिसकाव उसे मोटे तौर पर आधा कर
देते हैं)। उस खिड़की में कोई लाभकारी उत्परिवर्तन, जो सारे
उत्परिवर्तनों में $10^{3}$ से $10^{5}$ में एक है, किसी विराम-कूटक से
कहीं कम संभाव्य है, और आ भी जाए तो वह केवल $2s$ की प्रायिकता से स्थिर
होता है। इसलिए वह प्रति प्रायः तब तक मर चुकी होती है जब तक वरण उसके लिए
कोई काम ढूँढ़े --- और कशेरुकियों में द्विगुणित प्रतियों का देखा गया
अर्ध-आयु काल कुछ करोड़ वर्ष है।
\end{solution}

\begin{solution}{pb:b3:molecular-evolution:1}
\textbf{1.} $k = 0.012/(2\times 6.5\times 10^{6}) = 9.2\times 10^{-10}$
प्रति स्थल प्रति वर्ष; $\mu = 25k = 2.3\times 10^{-8}$ प्रति स्थल प्रति पीढ़ी।
\textbf{2.} $6\times 10^{9}\times 2.3\times 10^{-8} \approx 140$ नए
उत्परिवर्तन प्रति द्विगुणित जीनोम; उनका \qty{1.5}{\%}, यानी लगभग $2$,
कूटलेखी अनुक्रम में।
\textbf{3.} प्रति अगुणित जीनोम, $3\times 10^{9}\times 2.3\times 10^{-8}
\approx 70$ प्रतिस्थापन प्रति पीढ़ी स्थिर, जबकि हाल्डेन का
$1/300$: यानी उससे बीस हज़ार गुना अधिक जितना वरण चला सकता था, इसलिए
भारी बहुमत उदासीन है।
\textbf{4.} उस समष्टि में प्रति स्थल प्रति पीढ़ी
$2N\mu = 10^{5}\times 2.3\times 10^{-8} = 2.3\times 10^{-3}$
नए उत्परिवर्तन; उनका $1/2N =
10^{-5}$ स्थिर होगा।
\textbf{5.} $4N = \num{200000}$ पीढ़ियाँ, \qty{5}{Myr} --- यानी लगभग
उतना ही लंबा जितना उस विभाजन से बीता समय।
\textbf{6.} अपसरण उन उत्परिवर्तनों को गिनता है जो उनके साझा पूर्वज के
बाद उन दोनों अलग वंश-परंपराओं पर उठे; हर एक अपनी ही वंश-परंपरा में स्थिर
होता है, और लंबे चलन में वह दर $\mu$ है, चाहे हर एक जितना भी समय ले,
क्योंकि उत्परिवर्तन लगातार उठते हैं और वह देरी उस प्रवाह को टालती भर है,
बदलती नहीं। (इकलौता सुधार उस बहुरूपता के लिए है जो उस विभाजन पर मौजूद
थी, और जो दो युग्मविकल्पियों के साझा पूर्वज को उस जाति-विभाजन से पुराना
कर देती है।)
\textbf{7.} उदासीन $10^{-5}$; $s = 0.001$: $0.002$; $s = 0.01$: $0.0198$;
$s = 0.1$: $1 - \mathrm{e}^{-0.2} = 0.18$।
\textbf{8.} $s = 1/4N = 5\times 10^{-6}$: इससे नीचे अपवाह वरण को डुबो
देता है और वह उत्परिवर्तन उदासीन जैसा बरताव करता है; और इससे ऊपर वरण
उन संभावनाओं पर शासन करता है।
\textbf{9.} $s = -0.001$: $P = 0.002\,\mathrm{e}^{-200}$, यानी हर काम के
लिए शून्य --- उदासीन का $10^{-85}$। $s = -10^{-5}$, $4N|s| = 2$: $P =
2\times 10^{-5}/(\mathrm{e}^{2} - 1) = 3.1\times 10^{-6}$, यानी उदासीन का
\qty{31}{\%}: हलके हानिकारक उत्परिवर्तन स्थिर होते ही हैं।
\textbf{10.} प्रति स्थल प्रति पीढ़ी लाभकारी प्रतिस्थापन: $2N\times
10^{-5}\mu\times 2s = 10^{5}\times 10^{-5}\times 0.02\,\mu = 0.02\mu$;
उदासीन: $\approx\mu$। अनुकूली अंश $0.02/1.02 \approx \qty{2}{\%}$।
हालाँकि हर लाभकारी उत्परिवर्तन किसी उदासीन से दो हज़ार गुना अधिक स्थिर
होने योग्य है, वे इतने बिरले हैं कि पचास में केवल एक प्रतिस्थापन
अनुकूली है: वरण काम करता है और मुश्किल से दिखता है।
\textbf{11.} $p_{N} = 46/920 = 0.050$, $p_{S} = 84/280 = 0.30$। सुधरे हुए:
$d_{N} = -\tfrac{3}{4}\ln(1 - 0.0667) = 0.052$, $d_{S} = -\tfrac{3}{4}
\ln(1 - 0.40) = 0.38$। $\omega = 0.14$।
\textbf{12.} \omterm{def:b3:molecular-evolution:dnds}{शोधक वरण}, जो ऐमीनो-अम्ल बदलावों का लगभग \qty{86}{\%} हटा
देता है --- यानी कोई सामान्य जीन। $k = d_{S}/2T = 0.38/(1.8\times
10^{8}) = 2.1\times 10^{-9}$ प्रति स्थल प्रति वर्ष: यानी
मानव--चिंपैंज़ी मान से दुगुना, जैसा उस तुलना के लिए प्रत्याशित है जिसमें
तेज़ टिकती कृंतक वंश-परंपरा शामिल है; यानी वही परिमाण-कोटि।
\textbf{13.} $d = -\ln(1 - 0.55) = 0.80$ प्रति स्थल; $T = 0.80/(2\times
10^{-9}) = \qty{400}{Myr}$।
\textbf{14.} दो भिन्न शृंखलाएँ $\alpha_{2}\beta_{2}$ चतुष्क बनाती हैं,
जिसका सहकारी ऑक्सीजन-बंधन, \omterm{prop:b3:structural-biology:mechanism}{बोर प्रभाव} और 2,3-बिसफ़ॉस्फ़ोग्लिसरेट से
अन्यस्थानिक नियंत्रण असमान उपइकाइयों के बीच की उस अंतराफलक पर टिका है;
और फिर वह $\beta$ गुच्छा भिन्न बंधुताओं वाली भ्रूणीय, गर्भस्थ तथा वयस्क
शृंखलाओं में विविध हो सका।
\textbf{15.} $2kT = 2\times 8\times 10^{-9}\times 4\times 10^{7} = 0.64$
प्रतिस्थापन प्रति स्थल; और बीस ऐमीनो अम्लों के बीच संतृप्ति के साथ वह
कच्चा अंतर लगभग $0.95(1 - \mathrm{e}^{-0.64/0.95}) \approx 0.47$ है।
\qty{500}{Myr} पर $d = 8$: हर स्थल कई बार बदल चुका है और वह कच्चा अंतर
\qty{95}{\%} की अपनी छत पर बैठा है, और समय के बारे में कोई सूचना नहीं
ढोता।
\textbf{16.} अरब भर वर्षों में 100 स्थलों में दो वंश-परंपराओं के बीच
$10^{-11}\times 10^{9}\times 100\times 2 = 2$ प्रतिस्थापन। किसी
\qty{10}{Myr} विभाजन के लिए प्रत्याशा $0.02$ है: यानी लगभग हमेशा शून्य
अंतर, कोई विभेदन नहीं।
\textbf{17.} $1200\times 2\times 10^{-9}\times 0.04 = 9.6\times 10^{-8}$
प्रति वर्ष: पहला विराम लगभग \qty{10}{Myr} बाद प्रत्याशित है।
\textbf{18.} पूरे जीनोम को दुगुना करना हर जीन को एक साथ दुगुना कर देता
है, इसलिए अंतःक्रिया करते प्रोटीनों की रससमीकरणमिति बनी रहती है और कोई
मात्रा-असंतुलन उन प्रतियों के विरुद्ध वरण नहीं करता; जबकि कोई अकेली
अनुवर्ती प्रति किसी एक जीन की मात्रा बदल देती है और प्रायः हानिकारक
होती है, या तुरंत अनावश्यक और क्षय के लिए स्वतंत्र।
\textbf{19.} $\mathrm{e}^{-\num{80000}/\num{100000}} = \mathrm{e}^{-0.8} = 0.45$।
\textbf{20.} $\tfrac{2}{3}\times 0.45 = 0.30$: यानी देखे गए
\qty{30}{\%} से संगत।
\textbf{21.} आधा-आधा, यानी सारे वृक्षों का \qty{15}{\%} मनुष्य को
गोरिल्ला के साथ और \qty{15}{\%} चिंपैंज़ी को गोरिल्ला के साथ जोड़ता।
किसी एक की कोई स्पष्ट अधिकता उन दोनों जातियों के बीच उनके विभाजन के बाद
जीन-प्रवाह दिखाती, जो अपवाह पैदा नहीं कर सकता।
\textbf{22.} $N = \num{10000}$: $\mathrm{e}^{-4} = 0.018$, असहमति
\qty{1.2}{\%}। इसलिए देखा गया \qty{30}{\%} कोई पचास हज़ार की पुरखा
समष्टि माँगता है --- यानी आज के मनुष्यों के प्रभावी आकार से कई गुना।
\textbf{23.} जोड़ना यह मान लेता है कि हर जीन के पास वही \omterm{prop:b3:molecular-evolution:ils}{जाति-वृक्ष} है;
पर \omterm{prop:b3:molecular-evolution:ils}{अपूर्ण वंश-छँटाई} के साथ उन जीनों के पास अनेक वृक्ष हैं, और छोटी
भीतरी शाखाओं के लिए सबसे आम \omterm{prop:b3:molecular-evolution:ils}{जीन-वृक्ष} उस \omterm{prop:b3:molecular-evolution:ils}{जाति-वृक्ष} से भिन्न हो सकता
है, इसलिए अधिक आँकड़े उस ग़लत उत्तर को और आश्वस्त कर देते हैं। कोई
संगलन-विधि हर उम्मीदवार \omterm{prop:b3:molecular-evolution:ils}{जाति-वृक्ष} के लिए \omterm{prop:b3:molecular-evolution:ils}{जीन-वृक्षों} का प्रत्याशित
वितरण गणित करती है, और वह \omterm{prop:b3:molecular-evolution:ils}{जाति-वृक्ष} चुनती है जो देखी गई बारंबारताओं की
सबसे अच्छी व्याख्या करे।
\textbf{24.} \omterm{prop:b3:molecular-evolution:ils}{अपूर्ण वंश-छँटाई} सारे आधुनिक मनुष्यों की साझा पुरखा समष्टि
में हुई थी, इसलिए वह हर मानव समष्टि को नियंडरथलों से बराबर संबंधित बना
देती। ग़ैर-अफ़्रीकियों में कोई अधिकता का अर्थ है कि नियंडरथल
युग्मविकल्पी ग़ैर-अफ़्रीकियों के पूर्वजों के अफ़्रीका छोड़ने के बाद
आए: यानी मिश्रण, कोई पचास हज़ार वर्ष पहले।
\textbf{25.} $\mu \approx 2.3\times 10^{-8}$ प्रति स्थल प्रति पीढ़ी;
$\omega \approx 0.14$; और \omterm{prop:b3:molecular-evolution:ils}{जाति-वृक्ष} से असहमत \omterm{prop:b3:molecular-evolution:ils}{जीन-वृक्ष} लगभग
\qty{30}{\%}।
\end{solution}
